\section{Current mirror}
Phân tích mạch gương dòng Hình~\ref{fig:3-13} ($Q_1$, $Q_2$ giống nhau, $\beta = 100$), mô phỏng kiểm chứng; sau đó thay $R_1 = 100\,\Omega$, tính lại, mô phỏng và giải thích hiện tượng.

\given{fig3_13}{0.6\linewidth}{Current mirror circuit example\label{fig:3-13}}

\subsection{Tính toán lý thuyết ($R_1 = 2\,\text{k}\Omega$)}
\dotline{Theo định luật \hspace{3cm}, ta có:}
\dotline{$I_{CTRL} = $}
\dotline{$I_{C1} = $}
\dotline{$I_{B1} = I_{B2} = $}
\dotline{$I_L = $}

\subsection{Mô phỏng}
\simfigs{ex11_mirror_2k}{Mô phỏng gương dòng với $R_1 = 2\,\text{k}\Omega$}
\dotline{Vì sao gọi là mạch gương dòng (current mirror)?}
\dotlines{2}

\subsection{Thay $R_1 = 100\,\Omega$}
\dotline{$I_{CTRL} = $}
\dotline{$I_{C1} = $}
\dotline{$I_{B1} = I_{B2} = $}
\dotline{$I_L = $}

\simfigs{ex11_mirror_100}{Mô phỏng gương dòng với $R_1 = 100\,\Omega$}

\begin{notebox}{Hiện tượng \& giải thích}
\dotline{Hiện tượng:}
\dotlines{8}
\end{notebox}
